\documentclass[11pt]{article}
\usepackage[margin=1in]{geometry}
\usepackage{booktabs}
\usepackage{amsmath}
\usepackage{url}
\usepackage[hidelinks]{hyperref}
\usepackage{graphicx}

\title{Redundancy Is Not Duplication: Integrity Audits of Seven\\
Benchmark and Instruction Datasets}

\author{Ashish Sinha\\
\texttt{sinha.ashish.4.u@gmail.com}}

\date{}

\begin{document}
\maketitle

\begin{abstract}
Benchmark datasets are cited far more often than they are inspected. We audit
seven widely used datasets---Spider (train and dev), GSM8K (train and test),
HumanEval, BIRD-CRITIC 1.0, and Alpaca---for schema violations, degenerate
records, within-split duplication, and cross-split contamination, using a
single dependency-free procedure whose every number is reproducible.

We report two findings that bear on how such audits should be conducted. First,
the standard practice of clustering items by text similarity systematically
overcounts redundancy. In Spider's training split a similarity pass flags 77
duplicate clusters, but only 8 contain records identical across every evaluated
field; the remaining 69 reuse a question stem against a different database or
reference query, which is deliberate construction rather than redundancy. The
naive count overstates duplication by roughly $9\times$. The same pattern
appears in cross-split contamination: counting identical questions marks 6
Spider development items as contaminated, while additionally requiring the
reference SQL to match marks 2.

Second, measuring near-duplicate density across a sweep of similarity
thresholds, rather than at a single fixed threshold, exposes structure that a
point measurement hides. Human-authored datasets decay smoothly, gaining
$1.5\times$ to $2.5\times$ more clusters per threshold step. Alpaca is
essentially empty above $0.6$ and then gains $21.3\times$ in a single step. We
argue this discontinuity is a residue of the similarity filter applied during
Alpaca's generation, and that it has a direct practical consequence:
deduplicating a synthetic corpus at a threshold above its generation-time
filter returns zero and is mistaken for a clean result.

Code, per-dataset reports, and an interactive demonstration are available at
\url{https://github.com/ashishsinha1602/dataset-integrity-audit}.
\end{abstract}

\section{Introduction}

A benchmark score is only as meaningful as the distinctness of the items it is
computed over. Duplicate items inflate scores by letting a model bank the same
answer twice; overlap between a training split and an evaluation split inflates
them more directly still. Both are routinely assumed away.

The obvious way to check is to embed or shingle the natural-language field,
compute pairwise similarity, and count clusters above a threshold. We show that
this procedure, applied without qualification, is wrong in a specific and
correctable way, and that the correction changes the reported numbers by nearly
an order of magnitude.

Our contributions are:

\begin{enumerate}
\item A distinction between records that are \emph{identical across every
evaluated field} and records that merely \emph{share a text stem}, with
measurements showing the two are conflated by default and differ by $9\times$
on Spider and $4\times$ on BIRD-CRITIC.

\item The same correction applied to cross-split contamination: requiring the
reference answer to match, not only the question, reduces the apparent
contamination rate on Spider by $3\times$.

\item A threshold-sweep measurement (the \emph{duplicate curve}) that
distinguishes human-authored from filter-generated corpora, and reveals a
$21.3\times$ discontinuity in Alpaca consistent with a generation-time
similarity filter.

\item Audits of seven datasets, released with code that reproduces every
number.
\end{enumerate}

\section{Method}

\subsection{Field roles}

Datasets disagree on column names, sometimes inverting them: Spider's
\texttt{query} field holds SQL and \texttt{question} holds natural language,
which is the exact reverse of BIRD-CRITIC. We therefore parameterise every
measurement over three roles: \textbf{text} (the natural-language side),
\textbf{answer} (the reference solution), and \textbf{group} (the schema or
source an item is drawn from).

\subsection{Similarity}

Text is lowercased and whitespace-collapsed, then represented as the set of
$3$-word shingles. Similarity between items $i$ and $j$ is the Jaccard
coefficient $J(S_i, S_j) = |S_i \cap S_j| / |S_i \cup S_j|$. Reference answers
are compared exactly after stripping comments and collapsing whitespace;
near-identical SQL is common across genuinely independent problems and
contributes noise rather than signal.

Below $3{,}000$ records we compute all pairs exactly. Above it we use
MinHash-LSH ($128$ permutations, $32$ bands of $4$ rows) for candidate
generation only, and verify every candidate with its true Jaccard score, so
similarity values carry identical meaning under both paths. Shingles are hashed
with CRC-32 rather than Python's built-in \texttt{hash}, which is randomised
per process and would make bucketing irreproducible.

\subsection{Identical records versus shared stems}

Pairs above threshold are joined into clusters by union-find. Each cluster is
then classified: a cluster is \textbf{fully identical} if all members agree on
every evaluated field (text, answer, group, and any label fields), and a
\textbf{shared-stem variant} otherwise. Only the former is redundancy in the
sense that matters for scoring.

\subsection{Cross-split contamination}

For a training split $A$ and evaluation split $B$ we report four quantities:
items of $B$ whose text matches an item of $A$ exactly; items whose reference
answer matches; items where \emph{both} match the same record of $A$; and items
with a near-duplicate above threshold. We also report whether $A$ and $B$ share
any group values at all, since disjoint sources bound how much a text match can
mean.

\section{Results}

\subsection{Per-dataset audit}

Table~\ref{tab:main} reports the seven datasets. No dataset exhibits a single
schema violation. Alpaca carries 28 records with empty outputs; no other
dataset contains a degenerate record.

\begin{table}[h]
\centering
\small
\begin{tabular}{lrrrrrr}
\toprule
Dataset & Records & Schema & Degen. & Clusters & Identical & Variant \\
\midrule
Spider (train)        & 7{,}000  & 0 & 0  & 77 & 8 (16 rec., 0.23\%) & 69 \\
Spider (dev)          & 1{,}034  & 0 & 0  & 3  & 0 & 3 \\
GSM8K (train)         & 7{,}473  & 0 & 0  & 1  & 0 & 1 \\
GSM8K (test)          & 1{,}319  & 0 & 0  & 0  & 0 & 0 \\
HumanEval             & 164      & 0 & 0  & 0  & 0 & 0 \\
BIRD-CRITIC 1.0       & 600      & 0 & 0  & 4  & 1 (2 rec., 0.33\%)  & 3 \\
Alpaca                & 52{,}002 & 0 & 28 & 0  & 0 & 0 \\
\bottomrule
\end{tabular}
\caption{Integrity audit at Jaccard $\geq 0.70$. ``Identical'' counts clusters
whose members agree on every evaluated field; ``Variant'' counts clusters
sharing only a text stem.}
\label{tab:main}
\end{table}

The gap between the \emph{Clusters} and \emph{Identical} columns is the
finding. Spider's training split yields 77 clusters, of which 8 are genuine---a
$9.6\times$ overstatement. BIRD-CRITIC yields 4, of which 1 is genuine. In
BIRD-CRITIC the mechanism is explicit: instances \texttt{PostgreSQL\_215} and
\texttt{PostgreSQL\_216} pose an identical question filed under different
categories (\emph{Personalization} and \emph{Efficiency}) with different
reference SQL. A deduplication pass keyed on question text alone would discard
one of a deliberately constructed pair.

\subsection{Cross-split contamination}

Table~\ref{tab:leak} reports overlap between training and evaluation splits.

\begin{table}[h]
\centering
\small
\begin{tabular}{lrr}
\toprule
Check & Spider train $\to$ dev & GSM8K train $\to$ test \\
\midrule
Records compared              & 7{,}000 $\to$ 1{,}034 & 7{,}473 $\to$ 1{,}319 \\
Identical question            & 6 (0.58\%) & 0 \\
Identical reference answer    & 8 (0.77\%) & 0 \\
\textbf{Both identical}       & \textbf{2 (0.19\%)} & \textbf{0} \\
Near-duplicate question       & 7 (0.68\%) & 0 \\
Shared group values           & 0 databases & --- \\
\bottomrule
\end{tabular}
\caption{Cross-split contamination. Requiring both the question and the
reference answer to match is the measurement that bounds what a model can
answer from memory alone.}
\label{tab:leak}
\end{table}

Spider's splits share no databases, so an identical question across them is
posed of entirely different data. Every collision is a degenerate row-count
question landing on a coincidentally same-named table: \emph{``What is the
total number of airlines?''} appears against both \texttt{flight\_2} and
\texttt{flight\_4}, each reducing to \texttt{SELECT count(*) FROM AIRLINES}.
Two development items out of $1{,}034$ are answerable from memorised training
data, and both are trivial counts.

Reporting only the identical-question row would place the figure at 6 and the
framing at three times its true severity.

GSM8K shows no measurable overlap on any check. Because that result came from
the approximate path, we recomputed it under exact cross-product search over
all $9{,}852{,}887$ pairs: still zero, in 13 seconds. A zero produced by an
approximate method is precisely where verification is worth its cost.

\subsection{The duplicate curve}

Fixing a single threshold discards the shape of the similarity distribution.
Table~\ref{tab:curve} sweeps it.

\begin{table}[h]
\centering
\small
\begin{tabular}{lrrr}
\toprule
Threshold & Spider (train) & GSM8K (train) & Alpaca \\
\midrule
$\geq 0.7$ & 11.000 & 0.134 & 0.000 \\
$\geq 0.6$ & 20.571 & 0.401 & 0.019 \\
$\geq 0.5$ & 35.000 & 1.338 & 0.519 \\
$\geq 0.4$ & 52.714 & 3.345 & 11.076 \\
\midrule
$0.5 \to 0.4$ ratio & $1.5\times$ & $2.5\times$ & $\mathbf{21.3\times}$ \\
\bottomrule
\end{tabular}
\caption{Near-duplicate clusters per $1{,}000$ records across similarity
thresholds.}
\label{tab:curve}
\end{table}

The two human-authored datasets decay smoothly, gaining $1.5\times$ to
$2.5\times$ per threshold step, despite differing by two orders of magnitude in
absolute density. Alpaca is empty above $0.6$ and then gains $21.3\times$ in a
single step.

We interpret this discontinuity as a residue of the ROUGE-L similarity filter
applied when Alpaca's instructions were generated. We are explicit about what
is measured and what is inferred: the curve is measured, and the attribution to
the generation-time filter is an inference from its shape together with the
documented generation procedure. We did not have access to the generation
pipeline itself.

The filter was evidently effective. Alpaca contains \emph{no} exact duplicate
instruction across all $52{,}002$ records---verified by direct hashing,
independent of the approximate near-duplicate path---and none at all above
$0.6$ similarity.

The practical consequence is the point. A practitioner deduplicating a
synthetic corpus at a threshold at or above its generation-time filter recovers
zero and concludes the corpus is clean. The filter has already cleared that
range; the mass sits immediately beneath it. Sweeping the threshold rather than
fixing it makes this visible in a single pass.

\subsection{Agreement between exact and approximate search}

Switching search algorithms partway through a comparison is a standard way for
tables to quietly lose meaning. On Spider's development split, where both are
tractable, exact all-pairs search examined $534{,}061$ candidate pairs and
MinHash-LSH examined $232$. Both returned the same 3 pairs and the same 3
clusters.

\section{Limitations}

LSH offers no recall guarantee, so counts on datasets above the switching
threshold are lower bounds. The agreement check above found no loss, but a
single clean result is not a proof; every report we release states which method
produced it.

Similarity is computed over the text field only. The attribution of Alpaca's
curve discontinuity to its generation filter is an inference, as noted.

Cross-split contamination is measured only between the two splits supplied. It
says nothing about overlap with a model's pretraining corpus, which is the
larger and considerably harder question.

The $0.70$ threshold and $3$-word shingle size are defensible defaults rather
than tuned optima; both are exposed as parameters, and the curve measurement
exists precisely because no single threshold is privileged.

\section{Conclusion}

Across seven datasets and $69{,}592$ records, the six human-authored benchmarks
are structurally sound: no schema violations, no degenerate records, and 18
genuinely duplicated records in total. Reporting that plainly matters, because
the value of an audit lies in its willingness to return a negative result.

What is not sound is the default method for reaching such conclusions. Counting
similarity clusters overstates redundancy by $9\times$ on Spider; counting
identical questions overstates contamination by $3\times$. Both errors share a
correction---require the reference answer to match, not only the
question---and any pipeline that deduplicates on question text alone is
simultaneously discarding usable training data and overreporting contamination.

For synthetic corpora, the threshold sweep should replace the point
measurement. A single threshold cannot distinguish a corpus with no
near-duplicates from one whose near-duplicates were filtered out at generation
time and now sit just below wherever one happened to look.

\section*{Reproducibility}

All code is pure Python standard library; \texttt{datasets} is required only to
retrieve source data. Every table reproduces from the released repository:

\begin{verbatim}
python audit.py fetch --dataset xlangai/spider --split train --out spider.jsonl
python audit.py prep  --data spider.jsonl --preset spider
python audit.py leak  --a spider-train.jsonl --b spider-dev.jsonl --preset spider
python audit.py curve --data alpaca.jsonl --text-field instruction \
                      --answer-field output
\end{verbatim}

Repository: \url{https://github.com/ashishsinha1602/dataset-integrity-audit}\\
Interactive demonstration:
\url{https://huggingface.co/spaces/Ashsinha1/dataset-integrity-auditor}

\end{document}
